\subsection{QUOTIENT ACTION}

DEFINITION 4. Let \(\alpha \mapsto f_{\alpha}\) be an action of a set \(\Omega\) on a set \(\mathbf{E}\) . An equivalence relation \(\mathbf{R}\) on \(\mathbf{E}\) is said to be compatible with the given action if, for all elements \(x\) and \(y\) of \(\mathbf{E}\) such that \(x \equiv y \pmod{\mathbf{R}}\) and all \(\alpha \in \Omega\) , \(f_{\alpha}(x) \equiv f_{\alpha}(y) \pmod{\mathbf{R}}\) . The mapping which associates with an element \(\alpha \in \Omega\) the mapping of \(\mathbf{E} / \mathbf{R}\) into itself derived from \(f_{\alpha}\) by passing to the quotients is an action of \(\Omega\) on \(\mathbf{E} / \mathbf{R}\) called the quotient of the action of \(\Omega\) on \(\mathbf{E}\) .

Let E be a magma and R an equivalence relation on E. R is said to be left (resp. right) compatible with the law on E if it is compatible with the left (resp.

(right) action of E on itself derived from the law on E. For R to be compatible with the law on E it is necessary and sufficient that it be left and right compatible with the law on E.

We leave to the reader the statement and proof of the analogues of Propositions 6, 7 and 8 of § 1, no. 6.
% semantic-fragments: legacy-input-seam
\relax{}
